\chapter{Implementarea funcționalităților majore}

\section{Introducere}

Acest capitol detaliază implementarea principalelor funcționalități ale platformei
CourseHub: gestionarea conturilor și autentificarea, fluxul de completare a
profilului, căutarea și sistemul de cereri de colaborare, mesageria în timp real cu
partajare de fișiere, panoul de administrare și sistemul de localizare.

\section{Gestionarea utilizatorilor}

Entitatea \texttt{ApplicationUser} extinde clasa \texttt{IdentityUser} din
ASP.NET Identity cu proprietăți specifice domeniului:

\begin{lstlisting}[caption={Extinderea entității de utilizator (fragment)}]
public class ApplicationUser : IdentityUser
{
    public string? Bio { get; set; }
    public string? Role { get; set; }
    public string? EducationLevel { get; set; }
    public string? TeachingLevel { get; set; }
    public int? Age { get; set; }
    public string? Gender { get; set; }
    public string? SelectedTopics { get; set; }
    public string? AvatarPath { get; set; }
    public DateTime? BannedUntil { get; set; }
}
\end{lstlisting}

Câmpul \texttt{BannedUntil} ilustrează o decizie de proiectare deliberată: în loc
de un câmp boolean separat (\texttt{IsBanned}), o singură dată nullable acoperă
toate cazurile -- \texttt{null} înseamnă cont neblocat, o dată din viitor înseamnă
blocare temporară, iar \texttt{DateTime.MaxValue} este folosit convențional pentru
blocare permanentă. Verificarea „este blocat” se reduce, în orice context, la o
singură comparație: \texttt{BannedUntil > DateTime.UtcNow}.

\section{Autentificarea și verificarea prin OTP}

Înregistrarea unui cont nou nu autentifică imediat utilizatorul. După crearea
contului, se generează un cod numeric de 6 cifre, stocat împreună cu un termen de
expirare în tabela \texttt{OtpCode}, iar codul este trimis pe adresa de email
indicată. Contul rămâne neconfirmat (\texttt{EmailConfirmed = false}) până la
introducerea corectă a codului pe pagina dedicată de verificare. Același mecanism
este reutilizat pentru confirmarea modificărilor sensibile ale profilului
(schimbarea numelui de utilizator sau a adresei de email).

\section{Fluxul de completare a profilului}

Imediat după confirmarea contului, utilizatorul este direcționat către un flux de
completare a profilului, organizat în trei pași: alegerea rolului, completarea
detaliilor specifice rolului ales și selectarea materiilor de interes. O atenție
particulară a fost acordată transmiterii corecte a rolului între pași -- valoarea
aleasă la primul pas trebuie păstrată într-un câmp ascuns al formularului, altfel
serverul nu ar putea decide dacă trebuie salvat nivelul de educație sau nivelul de
predare. La finalul primului pas, rolul ales este scris atât în câmpul
\texttt{Role} al utilizatorului, cât și în tabela reală de roluri Identity:

\begin{lstlisting}[caption={Sincronizarea rolului domeniu cu rolul Identity}]
user.Role = Wizard.Role;
await _userManager.UpdateAsync(user);

var otherRole = Wizard.Role == "Teach" ? "Learn" : "Teach";
if (await _userManager.IsInRoleAsync(user, otherRole))
    await _userManager.RemoveFromRoleAsync(user, otherRole);
if (!await _userManager.IsInRoleAsync(user, Wizard.Role))
    await _userManager.AddToRoleAsync(user, Wizard.Role);
\end{lstlisting}

\section{Căutarea și sistemul de cereri de colaborare}

Căutarea permite filtrarea utilizatorilor după nume, rol, nivel de experiență
și materie, excluzând întotdeauna conturile cu rol de administrator din rezultate,
indiferent de valoarea câmpului \texttt{Role} pe care acestea o au. Un elev poate
trimite unei persoane cu rol de profesor o cerere de colaborare
(\texttt{HireRequest}); profesorul o poate accepta sau respinge din pagina
proprie de cereri, iar elevul poate vizualiza starea propriilor cereri trimise și
le poate anula cât timp sunt în așteptare. Odată o cerere acceptată, cei doi
utilizatori devin parteneri și pot iniția o conversație.

\section{Mesageria în timp real și partajarea fișierelor}

Comunicarea dintre utilizatori se realizează printr-un chat cu livrare
instantanee, implementat cu SignalR. Hub-ul expune o singură metodă, apelată de
client pentru a trimite un mesaj text:

\begin{lstlisting}[caption={ChatHub -- livrarea mesajelor în timp real}]
public class ChatHub : Hub
{
    public async Task SendMessage(
        string senderId, string receiverId, string message)
    {
        await Clients.User(receiverId)
            .SendAsync("ReceiveMessage", senderId, message);
    }
}
\end{lstlisting}

În paralel cu livrarea prin SignalR, clientul trimite mesajul și către un
endpoint API dedicat, responsabil de persistarea acestuia în baza de date --
separarea celor două căi permite ca istoricul conversației să rămână disponibil
chiar dacă unul dintre utilizatori nu este conectat în momentul trimiterii.

Profesorii pot, suplimentar, atașa fișiere conversației. Restricția este aplicată
exclusiv pe server, nu doar prin ascunderea butonului de atașare din interfață
pentru elevi: tipul fișierului este verificat printr-o listă albă de extensii,
dimensiunea este limitată, iar numele fișierului salvat pe disc este generat
aleator (\texttt{Guid.NewGuid}), pentru a preveni suprascrierea accidentală a
unui fișier existent cu același nume.

\section{Panoul de administrare}

Panoul de administrare este restricționat prin atributul
\texttt{[Authorize(Roles = "Admin")]}, aplicat la nivel de controller -- framework-ul
respinge orice cerere provenită de la un utilizator neautorizat, înainte ca vreo
linie de cod din interior să se execute.

Administratorul poate promova sau retrograda alți utilizatori, poate șterge un
cont (împreună cu toate mesajele, cererile de colaborare, codurile OTP și avatarul
asociate) și poate bloca temporar sau permanent un cont:

\begin{lstlisting}[caption={Aplicarea unei blocări (fragment din AdminController)}]
switch (duration)
{
    case "1d": user.BannedUntil = DateTime.UtcNow.AddDays(1); break;
    case "7d": user.BannedUntil = DateTime.UtcNow.AddDays(7); break;
    case "1m": user.BannedUntil = DateTime.UtcNow.AddMonths(1); break;
    case "permanent": user.BannedUntil = DateTime.MaxValue; break;
    case "unban": user.BannedUntil = null; break;
}
\end{lstlisting}

Un administrator nu se poate bloca sau șterge pe sine, iar niciun administrator
nu poate fi blocat sau șters de un alt administrator, pentru a preveni pierderea
accidentală a accesului la panoul de administrare. Pentru moderare, administratorul poate
vizualiza, pentru orice utilizator, lista completă a mesajelor trimise de acesta.

\section{Sistemul de localizare}

Textele din interfață sunt extrase în fișiere de resurse (\texttt{.resx}),
organizate în oglindă cu structura de directoare a vizualizărilor, iar textul
românesc este folosit direct drept cheie de resursă:

\begin{lstlisting}[caption={Exemplu de utilizare a localizării într-o vizualizare}]
<h1>@Localizer["Panou de administrare"]</h1>
\end{lstlisting}

Limba activă este determinată printr-un \emph{cookie} setat de un controller
dedicat, apelat la comutarea de limbă din bara de navigare; schimbarea limbii nu
necesită autentificare și este păstrată între sesiuni.
